\subsection{Introduction: The Generative Shift \quad / \quad \textthai{บทนำ: การเปลี่ยนแปลงเชิงสร้างสรรค์}}
\begin{paracol}{2}
Traditionally, physical science has been a "Forward Discipline." We define a structure—such as an atomic lattice or a mechanical bridge—and then use simulators (Density Functional Theory or Finite Element Analysis) to predict its properties. However, as the search space of possible materials grows into the quadrillions, forward simulation becomes a bottleneck.

This chapter introduces \textbf{Generative Modeling}\index{Generative Modeling@การสร้างแบบจำลองเชิงสร้างสรรค์ (Generative Modeling)}, the bridge to \textbf{Inverse Design}\index{Inverse Design@การออกแบบย้อนกลับ (Inverse Design)}. Instead of asking "What are the properties of this structure?", we ask "What structure has these properties?". By using Variational Autoencoders (VAEs) and Generative Adversarial Networks (GANs), we can train AI to "dream" of new physical configurations that satisfy specific engineering constraints.

We will explore how latent space representations allow physicists to navigate the continuous manifold of material possibilities, discovering new superconductors, high-strength alloys, and efficient heat sinks through purely data-driven imagination.

\switchcolumn

\begin{thai}
ตามธรรมเนียมแล้ว วิทยาศาสตร์กายภาพเป็น "สาขาวิชาเชิงรุก (Forward Discipline)" เรากำหนดโครงสร้าง เช่น โครงผลึกอะตอมหรือสะพานกลไก จากนั้นใช้เครื่องจำลอง (เช่น Density Functional Theory หรือ Finite Element Analysis) เพื่อทำนายคุณสมบัติของมัน อย่างไรก็ตาม เมื่อพื้นที่การค้นหาของวัสดุที่เป็นไปได้เพิ่มขึ้นเป็นสี่ล้านล้าน การจำลองเชิงรุกจึงกลายเป็นคอขวด

บทนี้จะแนะนำ \textbf{การสร้างแบบจำลองเชิงสร้างสรรค์ (Generative Modeling)} ซึ่งเป็นสะพานเชื่อมไปสู่ \textbf{การออกแบบย้อนกลับ (Inverse Design)} แทนที่จะถามว่า "โครงสร้างนี้มีคุณสมบัติอย่างไร" เราจะถามว่า "โครงสร้างใดมีคุณสมบัติเหล่านี้" ด้วยการใช้ Variational Autoencoders (VAEs) และ Generative Adversarial Networks (GANs)เราสามารถฝึก AI ให้ "ฝัน" ถึงโครงร่างทางกายภาพใหม่ๆ ที่ตรงตามข้อจำกัดด้านวิศวกรรมเฉพาะเจาะจงได้

เราจะสำรวจว่าการแสดงแทนพื้นที่แฝง (Latent space representations) ช่วยให้นักฟิสิกส์สำรวจแมนิโฟลด์ต่อเนื่องของความเป็นไปได้ของวัสดุได้อย่างไร โดยค้นพบคู่ตัวนำยิ่งยวดใหม่ โลหะผสมที่มีความแข็งแรงสูง และแผงระบายความร้อนที่มีประสิทธิภาพผ่านจินตนาการที่ขับเคลื่อนด้วยข้อมูลล้วนๆ
\end{thai}
\end{paracol}
